% ============================================================================
\chapter{系统调用}
\label{ch:syscall}
% Stage 12 — 待实现
% ============================================================================

\section{本章目标}

\begin{itemize}
  \item 理解系统调用的本质：用户态 → 内核态的受控入口
  \item 实现 \texttt{int 0x80} 系统调用门
  \item 第一批 syscall：\texttt{write}、\texttt{exit}、\texttt{brk}
  \item 用户态 syscall 封装库
\end{itemize}

\section{背景知识：为什么需要系统调用}

% TODO:
% - 用户态 (Ring 3) 不能直接访问硬件、修改页表
% - 系统调用是"合法的特权提升"通道
% - int 0x80 (Linux 传统) vs sysenter/syscall (现代)
% - 我们先用 int 0x80（简单，调试方便）

\subsection{系统调用约定}

% TODO:
% - EAX = syscall number
% - EBX, ECX, EDX, ESI, EDI = 参数
% - 返回值在 EAX

\section{IDT 配置}

% TODO: 在 idt_set_gate(0x80, ...) 中 DPL=3，允许 Ring 3 触发

\section{Syscall 分发}

% TODO: syscall_handler_c → switch(eax) → sys_write / sys_exit / sys_brk

\section{第一批系统调用}

% TODO:
% - sys_write(fd, buf, count): fd=1 输出到串口
% - sys_exit(code): 终止当前进程
% - sys_brk(addr): 调整用户堆边界

\section{用户态封装}

% TODO: 内联汇编封装 syscall(nr, arg1, arg2, ...)

\section{验证}

% TODO: 用户态程序调 write(1, "hello\n", 6) → 串口输出

\section{本章小结}

系统调用桥接了用户态和内核态。有了 write/exit，用户程序能输出信息和正常退出。
下一章实现进程管理（PCB、fork、exec）。
